\section{運用・保守設計}

システムの運用・保守・セキュリティ対策に関しての情報を以下に示す。

\subsection{運用}
\begin{itemize}
    \item \textbf{運用体制とコンテンツ管理} \\
    システムの運用は、システム管理者および大会主催者がそれぞれの権限に応じて行う。
    システム管理者はWebフレームワークが提供する管理者用インターフェース（Django管理画面等）を活用して
    マスタデータや利用状況を管理し、大会主催者はシステム上の機能を通じて大会の進行や参加状況を管理する。
    
    \item \textbf{不正利用・荒らしへの対応} \\
    利用規約に違反するプレイヤーや荒らし行為が確認された場合、
    管理者が対象ユーザーのアカウント状態（\texttt{is\_active}フラグ等）を無効化することで、
    迅速にサービスの利用を停止させる措置をとる。
    
    \item \textbf{サーバ・データベースの運用} \\
    稼働環境の標準モニタリング機能を利用し、サーバのCPU使用率やデータベースの接続数など、
    基本的な負荷状況を定期的に監視する。異常なアクセスや高負荷を検知した際は、
    速やかに管理者に通知される体制を整える。
\end{itemize}

\subsection{保守}
\begin{itemize}
    \item \textbf{データの永続化とステータス管理（大会情報のアーカイブ化）} \\
    参加者が大会終了後も過去の大会情報や対戦履歴を検索できるようにするため、
    終了した大会データはデータベースから物理削除しない運用とする。
    中止となった大会のデータも同様に保存する。
    大会データは「開催予定」「開催中」「終了済み」「中止」のステータスで論理的に区別して管理し、
    大会一覧画面等では開催状況による絞り込み機能を設けることで、混同を防ぐ。

    \item \textbf{契約・解約データの管理} \\
    サブスクリプション解約時は、契約状態を「解約済み」に更新し、
    問い合わせや売上管理に必要な一定期間は契約履歴を保持する。
    一方で、不要となったデータや長期保持が不要な情報は運用ルールに従って整理・削除し、
    万が一の情報漏洩時にも不要なデータが残らないよう管理する。
    また、大会契約のプラン変更時のデータ更新方針については、運用・実装要件に合わせて整合性を保つ。
    
    \item \textbf{定期メンテナンスとバックアップ} \\
    定期的にデータベースのバックアップを取得し、安全な領域に保管する。
    また、万が一の障害時に備え、バックアップデータからシステムが正常に復元できるかのテストを
    定期的に実施する。 
    
    \item \textbf{システムのアップデートと不具合修正} \\
    利用しているフレームワークやミドルウェアのサポート状況を監視し、
    必要に応じてセキュリティパッチの適用を行う。
    発見された不具合はバージョン管理システムを用いて修正履歴を管理し、
    動作確認の上で本番環境へ反映する。
    また、利用している外部パッケージの既知の脆弱性を、リリース前に検査し、以降も定期的に確認する。
\end{itemize}

\subsection{セキュリティとアクセス制御}

本番稼働中のシステムに対する不正アクセスや改ざん、情報漏洩を防ぐため、
Djangoが標準で備える防御機能と、本システムの仕様に合わせて追加実装・設定する対策を組み合わせて適用する。

\begin{itemize}
    \item \textbf{Django標準機能による対策} \\
    本番稼働中のWebアプリケーションに対する代表的な攻撃として、
    クロスサイトスクリプティング（XSS）はテンプレートの自動エスケープ、
    クロスサイトリクエストフォージェリ（CSRF）はCSRFミドルウェア、
    SQLインジェクションはORMによるパラメータ化クエリ、
    クリックジャッキングはX-Frame-Optionsヘッダにより対策する。
    パスワードはハッシュ化して保存し、平文では保持しない。

    \item \textbf{通信の保護} \\
    すべての通信をHTTPSとし、Nginxでの転送設定に加え、Django側でHSTS、
    セッションCookieおよびCSRF CookieのSecure属性、
    セッションCookieのHttpOnly属性とSameSite属性を設定する。

    \item \textbf{パスワード条件の検証} \\
    パスワードは12文字以上とし、文字種の指定は設けない。
    Django標準の検証機能により、最小文字数、よく使われるパスワードの拒否、
    氏名・メールアドレスとの類似の拒否、数字のみのパスワードの拒否を確認する。

    \item \textbf{パスワード再設定の安全性} \\
    再設定用リンクは、Django標準のトークンによる有効期限付き・一度限りのリンクとする。
    パスワードの更新後は、既存のログイン状態を終了する（3.2(e)と同一）。

    \item \textbf{データ単位の権限確認} \\
    一般ユーザー（競技者・観戦者）、大会主催者、システム管理者ごとにアクセス権限を付与する。
    利用者は1人1アカウントとし、同じ利用者が競技者にも主催者にもなれるため、
    権限は操作対象のデータごとに判定する。
    \begin{itemize}
        \item 大会の編集、競技者管理、トーナメント作成、結果の修正は、その大会の主催者のみ可能とする。
        \item 試合結果の入力と確認は、その試合に参加する競技者のみ可能とする。
        \item 当日受付は、参加登録済みの競技者本人のみ可能とする。
    \end{itemize}
    URLやIDの書き換えによって他人の大会や試合を操作できないよう、
    操作のたびにこの確認を行う。この確認はDjango標準では提供されないため、独自に実装する。

    \item \textbf{入力値の検証} \\
    大会および試合の配信URLは、httpまたはhttpsで始まるURLのみ登録できるようにし、
    リンクを経由した不正なスクリプトの実行を防ぐ。
    入力値の必須項目・形式・文字数は、サーバ側でも検証する。

    \item \textbf{管理画面・インフラ・機密情報の保護} \\
    Django管理画面は管理者のみが利用できるようにし、標準のURLから変更する。
    データベース（RDS）は外部に公開せず、セキュリティグループにより
    Web/APサーバ（EC2）からの接続のみを許可する。
    SECRET\_KEYやデータベースのパスワードなどの秘密情報は、
    ソースコードに記載せず環境変数で管理する。
    本番環境ではDEBUGを無効にし、ALLOWED\_HOSTSを設定する。
    また、決済情報等の機密データは外部決済サービスの仕組みを活用し、
    自システムで過剰に保持しないことで漏洩リスクを抑える。

    \item \textbf{ログイン試行回数の制限} \\
    パスワードの総当たり攻撃を防ぐため、同一のアカウントまたは接続元からのログイン失敗が
    一定回数に達した場合、一定時間ログインを受け付けないようにする。
    実装には、Django向けのライブラリ（django-axes等）を利用する。
\end{itemize}

\subsection{継続的な改善}

\begin{itemize}
    \item \textbf{フィードバックと機能追加} \\
    システムのリリース後は、利用者の意見を収集し、継続的な改善を行う。
    配信状況の可視化に対するニーズに応えるため、外部配信サービスとの連携機能などを
    開発状況や運用ニーズに応じて順次拡張していく計画とする。

    \item \textbf{セキュリティ対策の強化} \\
    リリース後は、運用状況に応じて、管理者アカウントへの二要素認証の導入や管理画面のIP制限、
    操作履歴・監査ログの拡充、WAFの導入、第三者による脆弱性診断などの対策を順次追加する。
\end{itemize}